%%%%%%%%%%%%%%%%%%%%%%%%%%%%%%%%%%%%%%%%%%%%%%%%%%%%%%%%%%%%%%%%%%%%%%%%%%%%%%%
% ERGÄNZUNGEN FÜR HAUPTPAPER
% Basierend auf Appendices U-AJ (440 Papers)
%%%%%%%%%%%%%%%%%%%%%%%%%%%%%%%%%%%%%%%%%%%%%%%%%%%%%%%%%%%%%%%%%%%%%%%%%%%%%%%

%==============================================================================
% ERGÄNZUNG FÜR KAPITEL 3: LIMITS OF UTILITY
% Neuer Abschnitt: Domain 3 - Temporal Preferences
%==============================================================================

\subsubsection{Domain 3: Temporal Preferences and Present Bias}

\paragraph{The Evidence.}
A vast literature documents systematic deviations from exponential discounting
\citep{frederick2002time}:
\begin{itemize}
  \item \textbf{Present bias:} People overweight immediate rewards relative to 
        future rewards, even controlling for pure time preference
  \item \textbf{Preference reversals:} The same person prefers \$110 in 31 days 
        over \$100 in 30 days, but \$100 today over \$110 tomorrow
  \item \textbf{Commitment demand:} People pay to restrict their future choices
  \item \textbf{Procrastination:} Costly tasks are perpetually postponed
\end{itemize}

\paragraph{The $\beta$-$\delta$ Model.}
\citet{laibson1997} formalized quasi-hyperbolic discounting:
\begin{equation}
  U_t = u(c_t) + \beta \sum_{\tau=1}^{T} \delta^\tau u(c_{t+\tau}) 
  \quad \text{where } \beta < 1
  \label{eq:beta-delta}
\end{equation}

The parameter $\beta$ captures present bias---the additional weight on immediate
consumption beyond standard exponential discounting $\delta$.

\paragraph{Empirical Estimates.}
Studies consistently find $\beta \approx 0.7$ and $\delta \approx 0.99$
\citep{frederick2002time,cohen2020measuring}. This implies:
\begin{itemize}
  \item Substantial present bias (30\% premium on immediate rewards)
  \item Modest long-run impatience
  \item Time-inconsistent behavior: plans made today are not followed tomorrow
\end{itemize}

\paragraph{Sophistication vs.\ Naïveté.}
\citet{PAP-odonoghue1999doing,odonoghue2001} distinguished:
\begin{itemize}
  \item \textbf{Sophisticates:} Know their $\beta < 1$, may seek commitment
  \item \textbf{Naïfs:} Believe $\hat{\beta} = 1$, perpetually surprised by own behavior
  \item \textbf{Partial naïfs:} $\beta < \hat{\beta} < 1$---underestimate but don't ignore bias
\end{itemize}

\paragraph{Framework Integration.}
Temporal preferences constitute the $\Psi_T$ dimension of context:
\begin{equation}
  \Psi_T = (\beta, \delta, \hat{\beta})
\end{equation}

Present bias implies that $C^*$ depends on the timing of decisions:
\begin{equation}
  C^*_{immediate} \neq C^*_{planned}
\end{equation}

This creates demand for commitment devices---institutional arrangements ($\Psi_I$)
that constrain future choices \citep{ashraf2006,bryan2010}. Examples include:
\begin{itemize}
  \item Illiquid savings accounts (401k, Christmas clubs)
  \item Deadlines and penalties
  \item Default enrollment in beneficial programs \citep{madrian2001}
\end{itemize}

\paragraph{Implications.}
Temporal context interacts with other $\Psi$ dimensions:
\begin{itemize}
  \item $\Psi_T \times \Psi_C$: Cognitive load worsens self-control
  \item $\Psi_T \times \Psi_I$: Institutions can provide commitment
  \item $\Psi_T \times \Psi_E$: Poverty may intensify present bias
\end{itemize}

%==============================================================================
% ERGÄNZUNG FÜR KAPITEL 3: LIMITS OF UTILITY
% Erweiterung Social Preferences mit Fehr-Charness (2024)
%==============================================================================

\paragraph{The State of the Art: Fehr \& Charness (2024).}
The comprehensive survey by \citet{fehrcharness2024} synthesizes three decades of
research on social preferences, documenting:
\begin{enumerate}
  \item \textbf{Distributional preferences:} Inequality aversion \citep{fehrschmidt1999},
        ERC \citep{boltonockenfels2000}, and social welfare concerns \citep{charnessrabin2002}
  \item \textbf{Belief-dependent preferences:} Reciprocity \citep{rabin1993} and
        guilt aversion \citep{charnessdufwenberg2006}
  \item \textbf{Image concerns:} Self-image and social image \citep{benaboutirole2006}
\end{enumerate}

\paragraph{Distribution of Types.}
A key finding: the majority of individuals have social preferences;
purely selfish individuals are a minority:
\begin{center}
\begin{tabular}{lcc}
\toprule
\textbf{Type} & \textbf{Students} & \textbf{General Population} \\
\midrule
Selfish & 40--60\% & 20--30\% \\
Inequality averse & 10--15\% & 15--20\% \\
Altruistic & 25--45\% & 45--55\% \\
\bottomrule
\end{tabular}
\end{center}

This heterogeneity is central to our framework: $\Psi_S$ varies across individuals,
and the distribution of types affects equilibrium outcomes.

\paragraph{Merit and Desert.}
Social preferences are context-dependent. Tolerance for inequality depends on
its source \citep{cappelen2007,almas2020}:
\begin{itemize}
  \item Inequality from effort/merit: More accepted
  \item Inequality from luck: Less accepted, more redistribution demanded
  \item Cultural variation: Americans more meritocratic, Scandinavians more egalitarian
\end{itemize}

%==============================================================================
% ERGÄNZUNG FÜR KAPITEL 5: COMPLEMENTARITY
% Integration Milgrom-Roberts
%==============================================================================

\subsection{The Milgrom-Roberts Framework}

The mathematical foundations of complementarity in economics were established by
\citet{milgromroberts1990compare,milgromroberts1990modern} and formalized by
\citet{topkis1998}.

\paragraph{Supermodularity.}
A function $f: X \to \mathbb{R}$ on a lattice $X$ is supermodular if:
\begin{equation}
  f(x \vee y) + f(x \wedge y) \geq f(x) + f(y)
  \label{eq:supermodular}
\end{equation}

For twice-differentiable functions, this is equivalent to:
\begin{equation}
  \frac{\partial^2 f}{\partial x_i \partial x_j} \geq 0 \quad \forall i \neq j
\end{equation}

This is precisely our complementarity condition $C_{ij} \geq 0$.

\paragraph{Consequences for Equilibrium.}
\citet{milgromroberts1990compare} showed that in supermodular games:
\begin{enumerate}
  \item Pure strategy Nash equilibria exist (no convexity needed)
  \item Equilibria form a complete lattice (smallest and largest exist)
  \item Comparative statics are monotone (higher parameters $\to$ higher equilibrium)
\end{enumerate}

\paragraph{Modern Manufacturing.}
\citet{milgromroberts1990modern} applied this to explain why modern manufacturing
practices cluster: JIT, TQM, flexible production, worker training, and flat hierarchies
are complements. Adopting one raises the return to adopting others.

In our notation:
\begin{equation}
  C_{JIT,TQM} > 0, \quad C_{JIT,flex} > 0, \quad C_{TQM,training} > 0, \ldots
\end{equation}

This explains why firms adopt coherent ``bundles'' of practices rather than
mixing traditional and modern approaches \citep{ichniowski1997}.

\paragraph{Strategic Fit.}
\citet{milgromroberts1995} extended to strategy: successful firms achieve
``fit'' across activities. Southwest Airlines succeeds through complementary
choices (point-to-point, no meals, fast turnaround, one aircraft type).
Imitating one element without others fails.

\paragraph{Framework Integration.}
Complementarity is not just a technical assumption---it is the core mechanism
through which context ($\Psi$) shapes equilibrium ($C^*$):
\begin{equation}
  \Psi \to C_{ij}(\Psi) \to C^*(\Psi)
\end{equation}

When $\Psi$ changes, complementarities change, and the entire equilibrium shifts---
not just marginal adjustments but potentially discontinuous transitions between
coherent configurations.

%==============================================================================
% ERGÄNZUNG FÜR KAPITEL 9: CONTEXT ENDOGENOUS
% Integration Spatial and Path Dependence
%==============================================================================

\subsection{Spatial Context and Agglomeration}

Physical location constitutes a fundamental dimension of context ($\Psi_P$).
The New Economic Geography \citep{krugman1991} showed how spatial equilibria
emerge from the interaction of increasing returns, transport costs, and
factor mobility.

\paragraph{Agglomeration Economies.}
\citet{marshall1890} identified three sources of agglomeration benefits:
\begin{enumerate}
  \item \textbf{Labor pooling:} Thick markets reduce matching frictions
  \item \textbf{Input sharing:} Specialized suppliers viable at scale
  \item \textbf{Knowledge spillovers:} Ideas flow locally
\end{enumerate}

In our notation, these are spatial complementarities:
\begin{equation}
  C_{ij}^{spatial} > 0 \quad \text{(co-location is complementary)}
\end{equation}

\paragraph{Urban Wage Premium.}
\citet{glaeser2008,glaeser2009} document that urban workers earn 25--35\% more
than rural workers with similar characteristics. Part is selection (skilled
workers move to cities), but substantial productivity gains from agglomeration
remain after controlling for sorting.

\citet{rosenthal2004,combes2015} estimate agglomeration elasticities:
\begin{equation}
  \frac{\partial \ln(\text{productivity})}{\partial \ln(\text{density})} \approx 0.03\text{--}0.08
\end{equation}

Doubling density raises productivity 3--8\%.

\paragraph{Spatial Path Dependence.}
\citet{davis2002} studied Japanese cities after WWII bombing. Despite massive
destruction, cities returned to pre-war size within 15--20 years. Geographic
fundamentals determine city location; history selects among multiple equilibria.

\citet{bleakley2012} found that historical portage sites (where goods were
carried around waterfalls) became cities that persist long after portage
became obsolete. Spatial $C^*$ exhibits strong path dependence:
\begin{equation}
  C^*_{spatial,t} = f(C^*_{spatial,t-1}, \Psi_P^{historical})
\end{equation}

\paragraph{Spatial Misallocation.}
\citet{hsiehmoretti2019} estimate that housing constraints in productive cities
(San Francisco, New York) reduce US GDP by 36\%. Workers cannot move to where
they would be most productive because housing supply is restricted.

This is spatial misallocation:
\begin{equation}
  Q^{actual} < Q^{optimal} \quad \text{due to spatial frictions}
\end{equation}

Policy implication: place-based policies \citep{klinemoretti2014} can be
efficient when multiple spatial equilibria exist and agglomeration externalities
are not internalized.

%==============================================================================
% ERGÄNZUNG FÜR KAPITEL 9: CONTEXT ENDOGENOUS
% Integration Complexity Economics
%==============================================================================

\subsection{Complexity and Emergence}

Standard economics assumes equilibrium exists, is unique, and is reached quickly.
Complexity economics \citep{arthur2015} challenges all three assumptions.

\paragraph{Economy as Complex Adaptive System.}
\citet{arthur1997} characterize the economy as:
\begin{itemize}
  \item \textbf{Heterogeneous agents:} Not identical, not representative
  \item \textbf{Local interaction:} Agents interact with neighbors, not all
  \item \textbf{Bounded rationality:} $\Psi_C < \infty$
  \item \textbf{Adaptation:} Agents learn, strategies evolve
  \item \textbf{Emergence:} Macro patterns not in micro rules
\end{itemize}

\paragraph{Path Dependence and Lock-In.}
\citet{arthur1989,arthur1994} showed that with increasing returns, small early
events can have large long-run consequences:
\begin{equation}
  C^*_{selected} \in \{C^*_1, C^*_2, \ldots\} \quad \text{history selects which}
\end{equation}

Multiple equilibria exist; path determines which one is reached. The selected
equilibrium may be inefficient---lock-in to inferior technology or institution.

\paragraph{Emergence.}
\citet{schelling1971,schelling1978} demonstrated that macro patterns can emerge
from micro behavior in unexpected ways:
\begin{equation}
  C^*_{macro} \neq \sum_i C^*_i \quad \text{(emergent properties)}
\end{equation}

His segregation model shows that mild individual preferences for similar neighbors
produce complete segregation---an outcome far exceeding individual preferences.

\paragraph{Framework Integration.}
Complexity economics provides the dynamic foundation for our framework:
\begin{itemize}
  \item $C^*$ is not calculated but \emph{emerges} from agent interaction
  \item Multiple $C^*$ may exist; history and $\Psi$ select among them
  \item Heterogeneity in $\Psi$ across agents is fundamental, not a nuisance
  \item Transitions between $C^*$ may be discontinuous (tipping points)
\end{itemize}

%==============================================================================
% ERGÄNZUNG FÜR KAPITEL 1: INTRODUCTION
% Stärkere Begründung für 8 Ψ-Dimensionen
%==============================================================================

\paragraph{Empirical Grounding of $\Psi$ Dimensions.}
Each dimension of $\Psi$ is derived from decision-theoretic requirements and
measured with established instruments:

\begin{itemize}
  \item \textbf{$\Psi_C$ (Cognitive):} Bounded rationality \citep{simon1962},
        prospect theory \citep{kahnemantversky1979}, dual-process theory
        \citep{kahneman2011}
  \item \textbf{$\Psi_I$ (Institutional):} Institutional economics \citep{PAP-north1990institutions},
        political institutions \citep{acemoglurobinson2012}
  \item \textbf{$\Psi_S$ (Social):} Social preferences \citep{fehrcharness2024},
        trust \citep{knackkeefer1997}
  \item \textbf{$\Psi_K$ (Informational):} Information economics
        \citep{akerlof1970,stiglitz2000}
  \item \textbf{$\Psi_T$ (Temporal):} Time preferences \citep{frederick2002time},
        present bias \citep{laibson1997}
  \item \textbf{$\Psi_P$ (Physical):} Economic geography \citep{krugman1991},
        urban economics \citep{glaeser2008}
  \item \textbf{$\Psi_E$ (Economic):} Resource constraints, budget effects
  \item \textbf{$\Psi_{Cult}$ (Cultural):} Cultural dimensions \citep{hofstede1980},
        world values \citep{inglehart2014}
\end{itemize}

The Global Preference Survey \citep{falk2018} provides cross-country measures
for patience ($\Psi_T$), risk tolerance, social preferences ($\Psi_S$), and
trust, demonstrating substantial and systematic variation across 76 countries.

%%%%%%%%%%%%%%%%%%%%%%%%%%%%%%%%%%%%%%%%%%%%%%%%%%%%%%%%%%%%%%%%%%%%%%%%%%%%%%%
\end{document}
